% Authority: German Wikiversity pageid 131353, revid 1111802, 2026-08-15T16:41:12Z.
% Exact UTF-8 witness SHA-256: b07aac01b6f803481fab9b5331e19480fb9bf0058597159b6152c65c7bc955ce.

\zwischenueberschrift{Lösung zu Aufgabe 1}

Es ist

\mavergleichskettedisp
{\vergleichskette
{ v
}
{ \in }{ T_PZ
}
{ = }{ \operatorname{kern} \left((D\varphi)_P\right)
}
{ } {
}
{ } {
}
}
{}{}{.}
Nach dem
\definitionsverweis {Satz über implizite Abbildungen}{}{}
gibt es eine offene Menge

\mathbed {P \in W} {}
 {W \subseteq G} {}
 {} {} {} {,}
eine offene Menge

\mavergleichskette
{\vergleichskette
{ V
}
{ \subseteq }{ \R^{n-m}
}
{ } {
}
{ } {
}
{ } {
}
}
{}{}{}
und eine stetig differenzierbare Abbildung
\maabbdisp {\psi} { V } { W } {}
derart, dass

\mavergleichskette
{\vergleichskette
{ \psi(V)
}
{ \subseteq }{ Z \cap W
}
{ } {
}
{ } {
}
{ } {
}
}
{}{}{}
ist und $\psi$ eine
\definitionsverweis {Bijektion}{}{}
\maabbdisp {\psi} { V } { Z \cap W } {}
induziert. Es sei

\mavergleichskette
{\vergleichskette
{ Q
}
{ \in }{ V
}
{ } {
}
{ } {
}
{ } {
}
}
{}{}{}
der Punkt mit

\mavergleichskettedisp
{\vergleichskette
{ \psi(Q)
}
{ =} { P
}
{ } {
}
{ } {
}
{ } {
}
}
{}{}{.}
Die Abbildung $\psi$ ist in jedem Punkt

\mavergleichskette
{\vergleichskette
{ Q
}
{ \in }{ V
}
{ } {
}
{ } {
}
{ } {
}
}
{}{}{}
\definitionsverweis {regulär}{}{}
und für das
\definitionsverweis {totale Differential}{}{}
von $\psi$ gilt

\mavergleichskettedisp
{\vergleichskette
{ (D\varphi)_P \circ (D\psi)_Q
}
{ =} { 0
}
{ } {
}
{ } {
}
{ } {
}
}
{}{}{,}
also

\mavergleichskettedisp
{\vergleichskette
{ \operatorname{bild} \left((D\psi)_Q\right)
}
{ =} { \operatorname{kern} \left((D\varphi)_P\right)
}
{ =} { T_PZ
}
{ } {
}
{ } {
}
}
{}{}{.}
Wegen der Regularität von $\psi$ in $Q$ ist
\maabbdisp {(D\psi)_Q} { \R^{n-m} } { \R^n } {}
injektiv und
\maabbdisp {(D\psi)_Q} { \R^{n-m} } { T_PZ } {}
bijektiv. Es sei

\mavergleichskette
{\vergleichskette
{ u
}
{ \in }{ \R^{n-m}
}
{ } {
}
{ } {
}
{ } {
}
}
{}{}{}
das Urbild von $v$ und sei
\maabbeledisp {\theta} { ]-\epsilon,\epsilon[ } { \R^{n-m} } {t} {Q+tu} {,}
wobei $\epsilon$ hinreichend klein gewählt sei, dass das Bild ganz in $V$ liegt. Dann besitzt
\maabbdisp {\gamma=\psi\circ\theta} { ]-\epsilon,\epsilon[ } { \R^n } {}
die Eigenschaft

\mavergleichskettedisp
{\vergleichskette
{ \gamma(0)
}
{ =} { \psi(\theta(0))
}
{ =} { \psi(Q)
}
{ =} { P
}
{ } {
}
}
{}{}{}
und

\mavergleichskettedisp
{\vergleichskette
{ \gamma'(0)
}
{ =} { (D\psi)_Q(\theta'(0))
}
{ =} { (D\psi)_Q(u)
}
{ =} { v
}
{ } {
}
}
{}{}{.}
